% Generated draft: requires mathematical and source review.
\documentclass[11pt]{article}
\usepackage[margin=25mm]{geometry}
\usepackage{fontspec}
\setmainfont{Latin Modern Roman}
\usepackage{amsmath,amssymb,mathtools,unicode-math}
\setmathfont{Latin Modern Math}
\usepackage{enumitem,needspace,xurl,hyperref}
\setlength{\parindent}{0pt}
\setlength{\parskip}{5pt}
\setlength{\emergencystretch}{3em}
\newcommand{\N}{\mathbb{N}}
\newcommand{\Z}{\mathbb{Z}}
\newcommand{\Q}{\mathbb{Q}}
\newcommand{\R}{\mathbb{R}}
\newcommand{\C}{\mathbb{C}}
\newcommand{\F}{\mathbb{F}}
\newcommand{\A}{\mathbb{A}}
\newcommand{\PP}{\mathbb{P}}
\newcommand{\E}{\mathbb{E}}
\newcommand{\eps}{\varepsilon}
\newcommand{\dd}{\,\mathrm{d}}
\begin{document}
\section*{UnsolvedMath ID: 700239}
\textbf{Source problem code:} catalogue entry
\subsection{Definitions and mathematical statement}
For a group \(G\), write \([x,y]=x^{-1}y^{-1}xy\). Let \(\mathcal N_2\) denote the class of all groups satisfying
\[
\forall x,y,z\in G,\qquad [[x,y],z]=1.
\]
Thus \(\mathcal N_2\) consists of groups of nilpotency class at most \(2\), including abelian groups and the trivial group; no finiteness restriction is imposed.

A group quasi-identity is a universally quantified implication
\[
\forall x_1,\ldots,x_n,\qquad
\left(\bigwedge_{i=1}^{m}u_i(x_1,\ldots,x_n)=v_i(x_1,\ldots,x_n)\right)
\Longrightarrow u(x_1,\ldots,x_n)=v(x_1,\ldots,x_n),
\]
where \(m\geq 0\) is finite and all displayed terms are group terms. A quasivariety of groups is the class of all groups satisfying some set of group quasi-identities.

Consider the lattice
\[
\mathsf L_q(\mathcal N_2)
=\{\mathcal Q:\mathcal Q\text{ is a quasivariety of groups and }\mathcal Q\subseteq\mathcal N_2\},
\]
ordered by inclusion. For any family \((\mathcal Q_i)_{i\in I}\) in this lattice, its meet is the intersection within \(\mathcal N_2\), and its join is
\[
\bigvee_{i\in I}\mathcal Q_i
=\bigcap\left\{\mathcal R\in\mathsf L_q(\mathcal N_2):
\bigcup_{i\in I}\mathcal Q_i\subseteq\mathcal R\right\}.
\]
The empty meet is \(\mathcal N_2\).

\textbf{Problem (M. I. Kargapolov, 1973).} Describe the lattice \(\mathsf L_q(\mathcal N_2)\), including its elements and their inclusion relations.

\textbf{Unresolved qualification.} The source says ``nilpotency class 2'' without explicitly distinguishing class exactly \(2\) from class at most \(2\). The formulation above uses class at most \(2\), the interpretation compatible with an ambient variety of groups. If exact class \(2\) was intended literally, clarification is necessary: every quasivariety of groups contains the trivial group, which does not have class exactly \(2\). The source also does not prescribe a particular form of lattice description.
\subsection{Short English statement}
Describe all quasivarieties contained in the class of groups of nilpotency class at most \(2\), and determine how they are ordered by inclusion. The phrase ``class 2'' requires the qualification explained above. The supplied registry labels this problem open; that status has not been independently verified.
\subsection{Sources}
\begin{itemize}
\item Local UnsolvedMath registry, permanent ID \(700239\), source code \texttt{catalogue entry}, title ``Kourovka Notebook Problem 4.31''; supplied complete source record, upstream revision \texttt{d371ee917ac9647fbb06cd6fc529deec6f3dda82}.
\item \emph{The Kourovka Notebook (Unsolved Problems in Group Theory)}, original problem \(4.31\), fourth issue (1973), proposed by M. I. Kargapolov. The supplied background identifies the twenty-first-issue edition (2026), \texttt{arXiv:1401.0300v48}; the supplied collection description names E. I. Khukhro and V. D. Mazurov as editors.
\item The local record supplies the status ``open'' and an update timestamp of \texttt{2026-10-16T14:00:00Z}. These are supplied registry metadata, not an independently verified literature assessment. No separate literature-verification date or problem-specific literature references are supplied.
\end{itemize}
\end{document}
